%------------------------------------------------------------------------------%
% Luis A. D'Afonseca
%
%------------------------------------------------------------------------------%
\documentclass[a4paper,12pt,fleqn]{article}

\usepackage{style_avaliacao}

% \ShowAnswers

\newcommand{\D}[2]{\frac{\partial {#1}}{\partial {#2}}}
\newcommand{\DS}[2]{\frac{\partial^2 {#1}}{\partial {#2}^2}}
\newcommand{\DM}[3]{\frac{\partial^2 {#1}}{\partial {#2} \partial {#3}}}

\newcommand{\vetor}[2]{\left(\begin{array}{c} #1 \\ #2 \end{array}\right)}

\newcommand{\veci}{\bm{i}}
\newcommand{\vecj}{\bm{j}}
\newcommand{\veck}{\bm{k}}

%------------------------------------------------------------------------------%
\begin{document}

\makeHeader{CFVVI}{Prova Substitutiva -- Turma 6}{11/09/2024}

\textbf{O exercício correspondente a prova que será substituida vale 50 pontos}

% 01
%------------------------------------------------------------------------------%
\exercise{25}
Calcule o limite
\(
  \lim_{(x, y) \to (0, 0)} \frac{xy}{x^2 + y^2}
\)
ou mostre que ele não existe.

\begin{answer}
  Restringindo a função à reta \( y = 0 \) temos
  \[
    \left(\lim_{(x, y) \to (0, 0)} \frac{xy}{x^2 + y^2}\right)  \evalat{y=0}{}
    = \lim_{x\to 0} \frac{0}{x^2}
    = \lim_{x\to 0} 0
    = 0
  \]
  Porém, se restringirmos $f$ à reta \( y=x\), temos
  \[
    \left(\lim_{(x, y) \to (0, 0)} \frac{xy}{x^2 + y^2}\right) \evalat{y=x}{}
    = \lim_{x\to 0} \frac{x^2}{x^2+x^2}
    = \lim_{x\to 0} \frac{1}{2}
    = \frac{1}{2}
  \]
  Os limites ao longo de diferentes caminhos são diferentes, logo o limite não existe.
  \clearpage
\end{answer}

% 02
%------------------------------------------------------------------------------%
\exercise{25}
Encontre e classifique os pontos críticos da função
\(
  f(x, y) = x^2 + y^2 - 4x - 6y + 12
\)

\begin{answer}
  Calculando as derivadas parciais de $f$
  \[
    f_x(x, y) = 2x - 4 \qquad\qquad f_y(x, y) = 2y - 6
  \]

  Como as derivadas parciais são polinômios elas existem em todos os pontos
  do plano. Portanto, temos apenas os pontos críticos onde as derivadas se anulam.
  Assim precisamos resolver o sistema \(\nabla f = 0\)
  \[
    \begin{cases}
      f_x(x, y) = 2x - 4 = 0 \\
      f_y(x, y) = 2y - 6 = 0
    \end{cases}
  \]
  que pode ser reescrito como
  \[
    \begin{cases}
      x = 2 \\
      y = 3
    \end{cases}
  \]
  Encontramos o ponto crítico \((2, 3)\)

  Para classificar o ponto crítico, usamos o teste da segunda derivada.

  Calculando as derivadas de segundas de $f$
  \[
    f_{xx}(x, y) = 2 \qquad\qquad
    f_{yy}(x, y) = 2 \qquad\qquad
    f_{xy}(x, y) = 0
  \]

  O discriminante é
  \[
    D(x, y) = f_{xx}(x, y)f_{yy}(x, y) - \big(f_{xy}(x, y)\big)^2
    = 2 \times 2 - 0^2
    = 4
  \]
  Como \(D(2, 3) = 4 >0\) o ponto é um máximo ou mínimo local,
  como \( f_xx(2, 3) = 2 > 0 \) o ponto é um mínimo local.
\clearpage
\end{answer}

% 03 ex 26 pg pdf 294
%------------------------------------------------------------------------------%
\exercise{25}
Encontre os números positivos $x$, $y$ e $z$, restritos a \(x+y+z^2=16\),
tais que seu produto seja máximo.

\begin{answer}
  Queremos maximizar
  \[
    f(x, y, z) = xyz
  \]
  com a restrição
  \[
    g(x, y, z) = x + y + z^2 = 16
  \]

  Para usar multiplicadores de Lagrange precisamos das derivadas parciais de $f$ e $g$
  \[
    f_x(x, y, z) = yz \qquad\qquad
    f_y(x, y, z) = xz \qquad\qquad
    f_z(x, y, z) = xy
  \]
  \[
    g_x(x, y, z) = 1 \qquad\qquad
    g_y(x, y, z) = 1 \qquad\qquad
    g_z(x, y, z) = 2z
  \]
  Vamos agora resolver o sistema \(\nabla f = \lambda\nabla g\) e \(g=16\)
  \[
    \begin{cases}
      yz = \lambda \\
      xz = \lambda \\
      xy = 2\lambda z \\
      x + y + z^2 = 16
    \end{cases}
  \]
  Igualando as duas primeiras equações
  \[
    yz = xz
  \]
  obtemos \(z=0\) ou \(y=x\).

  Como o enunciado especifica que os valores devem ser positivos então
  $z=0$ pode ser desconsiderado.

  Substituindo \(xz = \lambda\) e \(y=x\) na terceira equação
  \begin{align*}
    xy  & = 2\lambda z \\
    x^2 & = 2xz^2
  \end{align*}
  temos \(x=0\) ou \(z^2=\frac{x}{2}\).

  Como $x$ deve ser positivo temos
  \[
    y = x \qquad\text{e}\qquad z^2=\frac{x}{2}
  \]
  Substituindo na quarta equação
  \begin{align*}
    x + y + z^2 & = 16 \\[2mm]
    x + x + \frac{x}{2} & = 16 \\[2mm]
    \frac{2x}{2} + \frac{2x}{2} + \frac{x}{2} & = 16 \\[2mm]
    \frac{5x}{2} & = 16 \\[2mm]
    x & = \frac{16\times 2}{5} = \frac{32}{5}
  \end{align*}
  Portanto
  \[
    y = x = \frac{32}{5}
  \]
  \[
    z^2 = \frac{x}{2} = \frac{1}{2}\frac{32}{5} = \frac{16}{5}
    \qquad\qquad
    z = \pm \frac{4}{\sqrt{5}}
  \]
  Novamente, como $z$ deve ser positivo temos que os valores são
  \[
    x = y = \frac{32}{5} \qquad\qquad
    z = \frac{4}{\sqrt{5}}
  \]
\end{answer}

%------------------------------------------------------------------------------%
\end{document}
%------------------------------------------------------------------------------%
